\documentclass[11pt]{article}

\usepackage[margin=1in]{geometry}
\usepackage{amsmath,amssymb,amsthm}
\usepackage{graphicx}
\usepackage{enumitem}
\usepackage[hidelinks]{hyperref}

\newtheorem{theorem}{Theorem}
\newtheorem{corollary}[theorem]{Corollary}
\newtheorem{lemma}[theorem]{Lemma}
\newcommand{\dist}{\operatorname{dist}}

\title{A Uniform Geometric-Embedding Threshold\\
for Trees of Moderately Growing Degree}
\author{FA/Banach run \texttt{fa\_banach\_001}}
\date{9 August 2026}

\begin{document}
\maketitle

\section*{Status}

This packet records a candidate substantial partial answer to Problem~1.6 of
Altschuler--Dodos--Tikhomirov--Tyros \cite{ADTT}. The proof is complete modulo
the estimates already proved in that source, and the current verdict is
\[
 \boxed{\textbf{likely valid partial result; human review requested.}}
\]
It does not settle the full growing-degree problem or the optimal-constants
Problem~1.5.

\section*{Source problem}

A graph $G=(V,E)$ is \emph{geometrically embeddable} into a normed space $X$
if there is a map $\zeta:V\to X$ such that, for all distinct $u,v\in V$,
\[
 \|\zeta(u)-\zeta(v)\|_X\leq 1
 \quad\Longleftrightarrow\quad \{u,v\}\in E.
\]
The source proves a universal $\Theta(\log N/\log\log N)$ threshold when the
maximum degree $\Delta$ is fixed. On PDF page~3 it asks for the threshold when
$\Delta$ grows with $N$.

\begin{center}
\includegraphics[width=\textwidth]{figures/open_problem_crop.png}
\end{center}

In transcription, Problem~1.6 asks for an explicit function $m(N,\Delta)$ and
universal constants $\alpha_0,\alpha_1$ such that every $N$-vertex tree of
maximum degree $\Delta$ embeds into every normed space of dimension at least
$\alpha_1m$, while a complete tree of degree $\Delta$ does not embed below
$\alpha_0m$.

\section*{Proof Intuition}

The height of a complete degree-$\Delta$ tree with $N$ vertices is naturally
\[
 H=\frac{\log N}{\log(\Delta-1)}.
\]
Packing its mutually nonadjacent vertices forces dimension of order
$\log N/\log(2+H)$. This suggests the scale
\[
 q(N,\Delta)=\frac{\log N}{\log(2+H)}.
\]

The upper bound is hidden in the source's randomized construction. Its local
lemma separates paths into three ranges. Short paths are controlled with
probability $\exp(-m^{\kappa/2})$ and have dependency cost
$\exp(O(\log\Delta))$; this is the origin of the growth assumption below.
For paths of length at most $H$, the largest logarithmic dependency cost is
\[
 \frac{H}{\log H}\log\Delta\asymp \frac{\log N}{\log H}\asymp q.
\]
For paths longer than $H$, the probability exponent is $m\log H\asymp\log N$.
Thus $m=64q$ closes every local-lemma inequality. The Gaussian part of the
source proof contains no $\Delta$ and remains uniform on the whole interval
$c\log N/\log\log N\leq m\leq C\log N$.

\section*{The partial theorem}

Let $\kappa\in(0,1)$ be the absolute exponent in the thin-shell estimate
stated as Theorem~2.2 of \cite{ADTT}.

\begin{theorem}[Uniform moderately growing-degree threshold]\label{thm:main}
Let $(N_j)$ and $(\Delta_j)$ be integer sequences such that
$N_j\to\infty$ and $3\leq\Delta_j\leq N_j-1$. Set
\[
 D_j=\log(\Delta_j-1),\qquad
 H_j=\frac{\log N_j}{D_j},\qquad
 q_j=\frac{\log N_j}{\log(2+H_j)}.
\]
Assume
\begin{equation}\label{eq:growth}
 D_j=o\big(q_j^{\kappa/2}\big).
\end{equation}
Then, for all sufficiently large $j$, the following hold.
\begin{enumerate}[label=(\roman*)]
 \item Every tree on $N_j$ vertices of maximum degree at most $\Delta_j$
 geometrically embeds into every normed space of dimension at least
 $\lceil64q_j\rceil$.
 \item There is a breadth-first truncation of the complete rooted
 degree-$\Delta_j$ tree on $N_j$ vertices which does not geometrically embed
 into any normed space of dimension at most $q_j/2$.
\end{enumerate}
At cardinalities of a full complete rooted tree, the witness in (ii) is the
full complete tree. Hence $q(N,\Delta)$ is the threshold scale, up to universal
constants, throughout the regime \eqref{eq:growth}.
\end{theorem}

\begin{corollary}\label{cor:beta}
For every fixed $\beta<\kappa/2$, Theorem~\ref{thm:main} applies whenever
\[
 3\leq\Delta_N\leq \exp\big((\log N)^\beta\big).
\]
\end{corollary}

\section*{Proof}

We suppress the sequence index. Put $D=\log(\Delta-1)$,
$H=\log N/D$, and $q=\log N/\log(2+H)$.

\subsection*{Elementary consequences of the growth condition}

Since $D\geq\log2$, we have
\[
 H\leq \frac{\log N}{\log2},\qquad
 q\geq c\frac{\log N}{\log\log N},\qquad
 q\leq C\log N
\]
for absolute constants $c,C>0$. Conversely $q\leq C\log N$, so
\eqref{eq:growth} implies $D=o((\log N)^{\kappa/2})$. Therefore
$H=\log N/D\to\infty$, and
\begin{equation}\label{eq:logH}
 \frac{\log H}{\log(2+H)}\longrightarrow1.
\end{equation}

Retain the source parameters
\[
 c_1=\frac14,\qquad
 c_2=\min\left\{\frac14,\frac\kappa4\right\},\qquad
 \ell_0=\exp\big((\log N)^{c_2}\big),
\]
and its absolute integer $k_0$. Set
\[
 m=\lceil64q\rceil,\qquad k_1=\lfloor H\rfloor.
\]
Because $H\to\infty$ but $H=O(\log N)$, while $\ell_0$ grows faster than
every power of $\log N$, eventually
\begin{equation}\label{eq:order}
 k_0<k_1<\ell_0.
\end{equation}
Also
\begin{equation}\label{eq:smallterms}
 (\log N)^{c_2}=o(q^{\kappa/2})
 \quad\text{and}\quad
 (\log N)^{c_2}=o(q).
\end{equation}
Indeed, $c_2\leq\kappa/4$ and
$q\geq c\log N/\log\log N$.

\subsection*{Uniform local-lemma estimate}

Use exactly the random edge vectors, path events, dependency graph, and
weights from Section~4.2 of \cite{ADTT}. For a path-event of length $k$, the
weight is
\[
 b_k=
 \begin{cases}
 \displaystyle\frac{1}{2^{k+1}6\ell_0 k(\Delta-1)^k},&2\leq k\leq k_1,\\[6pt]
 \displaystyle\frac{3}{2\pi^2k^2N^2},&k_1<k\leq\ell_0.
 \end{cases}
\]
Claim~4.1 of the source bounds the number of dependent length-$k$ paths for a
fixed path by the minimum of $N^2$ and a constant multiple of
$\ell_0k(\Delta-1)^k$. If $k\leq k_1$, then
$(\Delta-1)^k\leq N$, while $k\ell_0=N^{o(1)}$. Hence, uniformly,
\begin{equation}\label{eq:dependency}
 6k\ell_0(\Delta-1)^k\leq N^2
\end{equation}
for all sufficiently large $N$. The product estimate (4.16) of the source
therefore remains valid:
\begin{equation}\label{eq:product}
 \prod_{B\sim A}(1-b_B)\geq\frac12
\end{equation}
for every path-event $A$. It remains to compare each event probability with
$b_k/2$.

For $2\leq k\leq k_0$, Lemma~3.3 of the source gives
\[
 \mathbb P(A)\leq \exp(-m^{\kappa/2}).
\]
On the other hand,
\[
 \log\frac{2}{b_k}leq C_0+k_0\log(2\Delta-2)+(\log N)^{c_2}
 =o(q^{\kappa/2})=o(m^{\kappa/2})
\]
by \eqref{eq:growth} and \eqref{eq:smallterms}. Thus
$\mathbb P(A)\leq b_k/2$ eventually.

Next suppose $k_0<k\leq k_1$. Lemma~3.2 of the source, with $c_1=1/4$, gives
\begin{equation}\label{eq:midprob}
 \mathbb P(A)\leq \exp\left(-\frac m8\log k\right).
\end{equation}
The function $x/\log x$ is increasing for $x\geq3$. Since
$\log(2\Delta-2)=D+\log2\leq2D$, \eqref{eq:logH} gives
\begin{align}
 \frac{k}{\log k}\log(2\Delta-2)
 &\leq \frac{k_1}{\log k_1}\log(2\Delta-2)\notag\\
 &\leq (2+o(1))\frac{HD}{\log H}
  =(2+o(1))\frac{\log N}{\log H}
  =(2+o(1))q.\label{eq:middlecost}
\end{align}
The constant, $\log k$, and $\log\ell_0$ terms in $\log(2/b_k)$ contribute
$o(q)\log k$ by \eqref{eq:smallterms}. Consequently
\[
 \log\frac{2}{b_k}\leq(3+o(1))q\log k,
\]
whereas $m/8\geq8q$. Comparing with \eqref{eq:midprob} yields
$\mathbb P(A)\leq b_k/2$ uniformly in this range.

Finally suppose $k_1<k\leq\ell_0$. The same small-ball estimate gives
\[
 \mathbb P(A)\leq
 \exp\left(-\frac m8\log k\right)
 \leq \exp\left(-8q\log H\right)
 =\exp\big(-(8-o(1))\log N\big),
\]
using \eqref{eq:logH}. Meanwhile
\[
 \log\frac{2}{b_k}=2\log N+2\log k+O(1)
 =(2+o(1))\log N
\]
because $k\leq\ell_0=N^{o(1)}$. Again $\mathbb P(A)\leq b_k/2$.

Together with \eqref{eq:product}, these three estimates verify the asymmetric
Lovasz local lemma \cite{AlonSpencer}. Thus the source's event $\mathcal L$,
on which every nonadjacent pair at tree distance at most $\ell_0$ has the
required uniform-vector separation, has positive probability.

\subsection*{Uniform Gaussian completion}

The remainder of Section~4.3 of \cite{ADTT} is independent of $\Delta$. We
record why every ``sufficiently large'' estimate there is uniform for the new
dimension. From the elementary bounds above,
\begin{equation}\label{eq:mrange}
 c\frac{\log N}{\log\log N}\leq m\leq C\log N.
\end{equation}
The source uses
\[
 \delta=\exp\left(-\left(\frac12-\frac{c_1}{2}\right)(\log N)^{c_2}\right).
\]
For adjacent pairs its Gaussian tail threshold is
$m^{-1/2}\delta^{-1}(\log N)^{-1}$. For distances $2\leq k\leq k_0$ it is,
up to absolute factors, $\delta^{-1}m^{-(3/2-\kappa)}$. For
$k_0<k\leq\ell_0$ it is, up to absolute factors,
$\exp((c_1/2)(\log N)^{c_2})m^{-1/2}$. In all three cases the exponential in
$(\log N)^{c_2}$ dominates every power of $m$ and $\log N$ allowed by
\eqref{eq:mrange}. The Gaussian upper-tail estimate therefore makes each bad
probability at most $N^{-2}$, exactly as in the source.

For $k>\ell_0$, the source's Gaussian small-ball bound is
\[
 \exp\big(-\tfrac{c_1}{4}m(\log N)^{c_2}\big).
\]
The lower bound in \eqref{eq:mrange} implies
$m(\log N)^{c_2}\gg\log N$, so this is also at most $N^{-2}$.

Conditioned on $\mathcal L$, fewer than $N^2/2$ unordered pairs remain. The
union bound is therefore strictly less than $1/2$, and a geometric embedding
exists with positive probability. If the ambient normed space has dimension
larger than $m$, perform the construction in any $m$-dimensional subspace.
This proves part~(i).

\subsection*{Complete-tree obstruction}

Use the breadth-first truncated complete tree from part~(A) of the source's
Theorem~1.1. It has $N$ vertices, maximum degree at most $\Delta$, a root $r$,
and radius
\[
 h_0=\left\lceil
 \frac{\log\left((N-1)(\Delta-2)/\Delta+1\right)}{\log(\Delta-1)}
 \right\rceil\leq H+1.
\]
At full complete-tree cardinalities, no truncation occurs.

Suppose it geometrically embeds into a $d$-dimensional normed space $X$.
Because the tree is bipartite, it has an independent set $V_0$ with
$|V_0|\geq N/2$. The radius-$1/2$ balls centered at the images of vertices in
$V_0$ are pairwise disjoint, and all lie in the radius-$(h_0+1/2)$ ball about
the image of $r$. Comparing Lebesgue volumes gives
\[
 \frac N2\leq(2h_0+1)^d.
\]
Hence
\begin{align*}
 d&\geq\frac{\log(N/2)}{\log(2h_0+1)}
 \geq\frac{\log(N/2)}{\log(2H+3)}\\
  &=(1-o(1))\frac{\log N}{\log(2+H)}=(1-o(1))q,
\end{align*}
where $H\to\infty$ was proved above. Thus $d\leq q/2$ is impossible for all
sufficiently large $N$, proving part~(ii).

For Corollary~\ref{cor:beta}, if
$D\leq(\log N)^\beta$ with $\beta<\kappa/2$, then
$H\geq(\log N)^{1-\beta}$ and
\[
 q\asymp\frac{\log N}{\log\log N}.
\]
It follows that $D=o(q^{\kappa/2})$, so Theorem~\ref{thm:main} applies.
\hfill$\Box$

\section*{Verification notes}

The verification audit checked every occurrence in Section~4 of the source
where $N$ was taken sufficiently large in terms of fixed $\Delta$.
The degree enters only through:
\begin{enumerate}
 \item the path-dependency count and inequality corresponding to
 \eqref{eq:dependency};
 \item the short-path comparison $m^{\kappa/2}\gg\log\Delta$;
 \item the transition length $k_1=\lfloor\log N/\log(\Delta-1)\rfloor$ in
 \eqref{eq:middlecost} and the longer-path estimate.
\end{enumerate}
All later estimates are degree-free, and \eqref{eq:mrange} is exactly the
uniform range they require. No finite computation is used as proof.

The main human-review focus should be the uniformity of the source's
Lemma~3.3 in $k\leq k_0$ (the lemma is stated uniformly for all $k\geq2$) and
the interpretation of ``complete tree'' at non-full cardinalities. The lower
bound is exact for full complete trees and uses the same truncation as the
source for arbitrary $N$.

\section*{Limitations and novelty status}

The theorem leaves open degrees for which
$\log(\Delta-1)$ is comparable to or larger than $q^{\kappa/2}$, does not
optimize the constants, and does not address Problem~1.5. Its use of the
thin-shell exponent means the displayed degree range improves automatically
with any stronger uniform short-path estimate.

A bounded literature search on 9 August 2026 used arXiv:2504.15212, the exact
paper title, the exact phrase ``Trees of growing maximal degree'', and keyword
combinations involving geometric embeddings, trees, growing degree, and
normed spaces. It found no separate paper answering Problem~1.6 or stating
this uniformized range. Since the source v2 dates from 17 April 2026, novelty
confidence is moderate rather than definitive.

\section*{Human-review recommendation}

Review the three local-lemma comparisons and the Gaussian uniformity interval
\eqref{eq:mrange}. If those checks pass, this packet gives a rigorous
substantial partial answer to Problem~1.6 with an explicit threshold scale and
universal constants on a nontrivial growing-degree regime.

\begin{thebibliography}{9}
\bibitem{ADTT}
D.~J. Altschuler, P.~Dodos, K.~Tikhomirov, and K.~Tyros,
\emph{A universal threshold for geometric embeddings of trees},
arXiv:2504.15212v2 (17 April 2026),
\url{https://arxiv.org/abs/2504.15212}.

\bibitem{AlonSpencer}
N.~Alon and J.~H. Spencer,
\emph{The Probabilistic Method}, 4th ed., Wiley, 2016.
\end{thebibliography}

\end{document}
